% Draft for issue #26. All claims and notation remain subject to independent review.
\section{Finite lattice gauge theory}
\label{sec:ex-lattice}
\label{ex-lattice}

The starting point is a probability measure on finitely many group-valued links.
Compactness makes this measure elementary to define; the difficulty of the
continuum problem is to control its observables uniformly as the number of links
grows. We first fix the action and its normalization.

\subsection{The periodic lattice and its probability law}
\label{subsec:ex-lattice-finite}

Fix a closed subgroup $G\subset U(r_G)$ in a specified faithful unitary
representation, with nonabelian simple real Lie algebra as in
Section~\ref{ex-statement}. The group is compact. For a positive
integer $\ell$, let
\[
  \Lambda_\ell=(\mathbb Z/\ell\mathbb Z)^4,\qquad
  E_\ell=\Lambda_\ell\times\{0,1,2,3\}.
\]
A link configuration is a function $U:E_\ell\to G$. The four coordinates are
Euclidean; this finite model does not distinguish a time direction in its
action. An oriented link in the negative direction is represented by the
inverse of the corresponding positively oriented link.

For $x\in\Lambda_\ell$ and distinct axes $i,j$, the positively oriented
plaquette holonomy based at $x$ is
\[
 U_{ij}(x)=U_i(x)U_j(x+e_i)U_i(x+e_j)^{-1}U_j(x)^{-1}.
\]
Count each pair $i<j$ once. The action is
\[
 S_\ell(U)=\sum_{x\in\Lambda_\ell}\sum_{0\leq i<j<4}
   \left(1-\frac{\operatorname{Re}\operatorname{tr}U_{ij}(x)}{r_G}\right).
 \label{eq:ex-lattice-action}
\]
Each summand is real and nonnegative: a unitary matrix has each diagonal
entry of modulus at most one, so $\operatorname{Re}\operatorname{tr}V\leq r_G$.
The convention has no extra factor of $1/2$ over faces. This action is
dimensionless. The bare
coefficient $\beta$ multiplies it in the Boltzmann weight; it is not the
blocking-kernel penalty used later in the RG construction.

Let $H_G$ be probability Haar measure on $G$, and let
\[
 \rho_\ell=\bigotimes_{(x,i)\in E_\ell}H_G,
 \quad Z_\ell(\beta)=\int e^{-\beta S_\ell(U)}\,d\rho_\ell(U),
 \quad d\mu_{\ell,\beta}(U)=Z_\ell(\beta)^{-1}e^{-\beta S_\ell(U)}d\rho_\ell(U).
\]
The partition function is finite and strictly positive. Indeed the link space
is a finite product of compact groups, hence compact, and $S_\ell$ is continuous;
therefore $e^{-\beta S_\ell}$ is bounded and integrable. It is strictly
positive at every configuration, so its integral against the probability
measure $\rho_\ell$ is positive. Integrating the displayed density gives
$Z_\ell(\beta)/Z_\ell(\beta)=1$, so $\mu_{\ell,\beta}$ is a probability
measure. This argument works for every real $\beta$ for which the displayed
definition is used; the construction theorem later selects its own coefficient
sequence.

\begin{proposition}[Gauge invariance of the finite action and law]
\label{prop:ex-lattice-gauge}
For any map $g:\Lambda_\ell\to G$, transform links by
\[
 (g\cdot U)_i(x)=g(x)U_i(x)g(x+e_i)^{-1}.
\]
Then $S_\ell(g\cdot U)=S_\ell(U)$ and $\mu_{\ell,\beta}$ is invariant under
$U\mapsto g\cdot U$.
\end{proposition}
\begin{proof}
Substitution into the four factors of $U_{ij}(x)$ cancels the group elements at
the three intermediate vertices. The remaining factors give
\[
 (g\cdot U)_{ij}(x)=g(x)U_{ij}(x)g(x)^{-1}.
\]
Trace is invariant under conjugation, since for matrices $A,B$ of matching
finite size $\operatorname{tr}(AB)=\operatorname{tr}(BA)$; hence every
plaquette summand and their finite sum are unchanged.

For each link $(x,i)$, the transformed value is obtained from $U_i(x)$ by left
and right multiplication by fixed group elements. Left and right translations
preserve Haar measure. The full transformation acts separately on the
coordinates of the finite product $E_\ell$, so the product Haar measure
$\rho_\ell$ is preserved. Since the density $Z_\ell(\beta)^{-1}
e^{-\beta S_\ell(U)}$ is invariant, the weighted measure is preserved as well.
This proves invariance of the law, not gauge invariance of every observable.
An insertion must separately satisfy the relevant invariance condition.
\end{proof}

\subsection{Reflection positivity on an even torus}
\label{subsec:ex-lattice-finite-reflection}

We next prove the finite-volume positivity needed for the reflection
construction. Fix an axis as the time coordinate, let $\ell=2h$ with
$h\geq1$, and take $\beta\geq0$. There are two cuts: the site reflection
$t\mapsto-t$, whose strict positive sector is $1\leq t<h$, and the link
reflection $t\mapsto-1-t$, whose positive sector is $0\leq t<h$; both
coordinates are taken modulo $2h$. Reflections based at other time
coordinates are obtained by translation. The reflected configuration pulls
back links along the reflected oriented edges; a reflected negative edge is
read as the inverse of its positive representative. Write $\Theta U$ for this
configuration.

\begin{proposition}[Finite-volume reflection positivity]
\label{prop:ex-lattice-finite-reflection}
For either cut, let $F$ be a bounded Borel, gauge-invariant function of the
link configuration which depends only on links with both endpoints in the
strict positive sector. Then
\[
 \int F(U)\,\overline{F(\Theta U)}\,d\mu_{2h,\beta}(U)\geq0.
\]
The inequality is real and the same statement holds at every translated cut
and for every choice of time axis. It applies at each finite volume; no
continuum limit is asserted.
\end{proposition}
\begin{proof}
We first fix the transfer-kernel normalization. Let $X$ and $Y$ be the
spatial link configurations on two adjacent time slices, let $u_v$ be the
intervening time link at spatial vertex $v$, and let $e=(v,v+e_j)$ run over
the positive spatial links in one slice. The temporal plaquette has holonomy
$X_eu_{v+e_j}Y_e^{-1}u_v^{-1}$. Put
\[
 H(X,Y)=\int\exp\!\left[-\beta\sum_e
 \left(1-\frac{\operatorname{Re}\operatorname{tr}
 (X_eu_{v+e_j}Y_e^{-1}u_v^{-1})}{r_G}\right)\right]
 \prod_v dH_G(u_v).
\]
This is the single common integration over the family of time links, not an
independent gauge average for each spatial edge. We claim that $H$ is a
positive-semidefinite kernel. Define the spatial gauge action
$(u\cdot Y)_e=u_vY_eu_{v+e_j}^{-1}$ and the real inner product
$\langle X,Y\rangle=\sum_e\operatorname{Re}\operatorname{tr}(X_eY_e^*)$.
For $c=\beta/r_G\geq0$, set
\[
 H_0(X,Y)=e^{-\beta |E_{\rm sp}|}\exp(c\langle X,Y\rangle).
\]
For any finitely many $X^a$ and complex coefficients $z_a$,
$\sum_{a,b}\overline{z_a}z_b\langle X^a,X^b\rangle^m\geq0$:
view the matrix entries as real coordinates and take the squared norm of
$\sum_a z_a(X^a)^{\otimes m}$ in the complexified real tensor space. The
exponential series, with nonnegative coefficients, therefore makes $H_0$
positive semidefinite. It is invariant under a simultaneous spatial gauge
transformation of both arguments. Product Haar invariance gives
\[
 H(X,Y)=\int H_0(X,u\cdot Y)\,du
       =\int\!\!\int H_0(v\cdot X,u\cdot Y)\,du\,dv.
\]
More explicitly, if $M_m(X)=\int(u\cdot X)^{\otimes m}\,du$, then
\[
 H(X,Y)=e^{-\beta |E_{\rm sp}|}\sum_{m\geq0}
 \frac{c^m}{m!}\langle M_m(X),M_m(Y)\rangle.
\]
Thus $H$ is real, symmetric, and positive semidefinite. The series converges
uniformly because the spatial-link space is compact. Products of such kernels
are positive semidefinite by taking tensor products of their feature spaces;
the same remains true when the two boundary slice variables coincide.

For a link cut, partition the temporal links into those wholly in the positive
half, those in its reflected half, and the two sets of crossing links. The
spatial links are partitioned between the two halves. Every spatial plaquette
belongs to one half, every non-crossing temporal plaquette belongs to one
half, and the remaining temporal plaquettes occur at the two cut boundaries.
Reflection identifies the two half actions. After integrating the crossing
links, the unnormalized reflected expectation is a double integral over the
two open half-configurations, with integrand
\[
 \Phi(x)\,\overline{\Phi(y)}\,
 H(X_1(x),X_1(y))\,H(X_2(x),X_2(y)).
\]
Here $\Phi$ is $F$ times the positive-half Boltzmann factor, and $X_1,X_2$
are its two boundary spatial slices. Insert the two uniformly convergent
feature expansions. Fubini turns each term into the squared norm of the
integral of $\Phi$ against the corresponding tensor feature, with a
nonnegative coefficient. Boundedness of $F$ and compactness justify the
interchange of integrals and sums. The integral is therefore nonnegative.
This also covers $h=1$, when the two boundary slice variables within one half
are the same.

For a site cut, condition instead on the spatial links in the two fixed
reflection planes. The remaining links occur in reflected pairs, and the
plaquette action splits into the action in one half, its reflection, and the
spatial plaquettes contained in the fixed planes. Since $F$ reads only links
with both endpoints in the strict positive sector, the two half integrations
are conjugates. Their product is an absolute square, multiplied by the positive
Boltzmann weight of the fixed-plane plaquettes. The same factorization works
when the strict sector is empty. Product Haar invariance supplies the reflected
change of variables, including inversion of reversed links. All integrands are
bounded and measurable on a compact finite product, so Fubini's theorem
applies. Finally $Z_{2h}(\beta)>0$; division by it preserves nonnegativity.
Translation of the cut and permutation of axes preserve the product prior and
plaquette action, proving the remaining cases.
\end{proof}

\subsection{Scales, observables, and the ordered limit}
\label{subsec:ex-lattice-scales}

For mesh index $N$ and physical-period index $p$, the family uses
strictly increasing positive-integer sequences $D_N$ and $P_p$. Set
\[
 a_N=D_N^{-1},\qquad \ell_{N,p}=P_pD_N=P_p/a_N,
 \qquad \Lambda_{N,p}=(\mathbb Z/\ell_{N,p}\mathbb Z)^4.
\]
Thus $\ell_{N,p}$ is a site count per axis, while $P_p$ is the physical linear
period. In the selected construction, every $P_p$ is divisible by four and
every $\beta_{N,p}$ is positive. Consequently Proposition~\ref{prop:ex-lattice-finite-reflection}
applies on every one of its finite $(N,p)$ grids. These extra properties are
of the constructed schedule and coupling; strict increase alone in the abstract
family definition does not impose either one. A site is embedded in Euclidean spacetime by multiplying the symmetric
least-absolute residue in each coordinate by $a_N$. A token $v$ specifies an
insertion $O_{N,p}(v,z;U)$. We require the insertion to be continuous and gauge invariant, with the
fixed-stencil locality specified in Section~\ref{ex-statement}.

The finite $m$-point moment is the integral of the product of its $m$
insertions under $\mu_{\ell_{N,p},\beta_{N,p}}$. Smearing uses a lattice sum
weighted by $a_N^{4m}$. Continuity and compactness of the finite configuration space give
integrability of each finite product. The
public convergence interface first demands, at each fixed
suitable physical volume, a full $N\to\infty$ mesh limit of these smeared
moments, and only then a volume limit along a specified strictly increasing
subsequence. It is not an arbitrary joint or diagonal limit. The support of a
test function must eventually fit inside the period box and its arguments must
remain separated as specified in Section~\ref{ex-statement}.

The next task is to control these moments as the number of links tends to
infinity. The required estimates must concern the same bare coefficients and
local insertions at all cutoffs and periods. The extraction argument in the
next section isolates a local modulus of continuity and a local size bound;
the renormalization construction must supply those estimates with the indicated
order of quantifiers.
